% Seminar 3 - Page 6. Fragment for
% % Seminar 3 - Page 6. Fragment for
% % Seminar 3 - Page 6. Fragment for
% % Seminar 3 - Page 6. Fragment for
% \input{seminars/pages/seminar3-pagina6-en}; no preamble or forced page break.
% EDITORIAL NOTES:
% - This page continues examples (i), (iii), and (iv) of the preceding page.
% - In the opening line, square brackets and braces denote the integer and
%   fractional parts; the compact handwritten expression is interpreted as
%   [x]+{x}-[y]-{y} being an integer when x~y.
% - The proof explicitly shows reflexivity and symmetry of congruence mod Z;
%   transitivity is not written out on this manuscript page.
% - A small marginal reminder repeats (iv), x~y iff x-y in Q, and R/~ ~ R.
%   It is not reproduced as another numbered example.
% - The diagrams are redrawn for legibility; p restricted to a chosen set
%   of representatives is the indicated bijection.
% - Equality of complex arguments depends on a consistent choice of argument
%   (or equality modulo 2pi); the notes leave this convention implicit.
% - In the last diagram the handwritten notation for the map to S' is faint;
%   the unit-normalisation map z -> z/|z| is consistent with the proof.

\subsection*{Seminar 3 - Page 6}
\begingroup
\setlength{\abovedisplayskip}{5pt}
\setlength{\belowdisplayskip}{5pt}
\setlength{\abovedisplayshortskip}{3pt}
\setlength{\belowdisplayshortskip}{3pt}

\noindent\textbf{Proof (continuation of example (i)).}
Write $x=\lfloor x\rfloor+\{x\}$ and $y=\lfloor y\rfloor+\{y\}$.
If $x-y\in\mathbb Z$, then
\[
 \lfloor x\rfloor+\{x\}-\lfloor y\rfloor-\{y\}\in\mathbb Z,
 \qquad \{x\}-\{y\}\in\mathbb Z.
\]
For the relation $x\sim y\iff x-y\in\mathbb Z$:
\[
 \begin{aligned}
 \text{(R)}\quad &x\sim x &&(x-x=0\in\mathbb Z),\\
 \text{(S)}\quad &x\sim y\Longrightarrow y\sim x
     &&(x-y\in\mathbb Z\Longrightarrow y-x\in\mathbb Z).
 \end{aligned}
\]

\noindent\underline{Variant 1.} The interval $[0,1)$ is a system of
representatives for $\sim$ if and only if
\[
 p|_{[0,1)}:[0,1)\longrightarrow\mathbb R/{\sim}
 \quad\text{is bijective}.
\]
\noindent\textbf{Step 1.} If $x,y\in[0,1)$ and $x\ne y$, then
$x\not\sim y$ (evident).

\noindent\textbf{Step 2.} For every $x\in\mathbb R$ there is
$t\in[0,1)$ with $x-t\in\mathbb Z$: take $t=\{x\}$.

\medskip
\noindent\textbf{(iii) Complex numbers with equal modulus.}
For the quotient $\mathbb C/{\sim}$, where
$z_1\sim z_2\iff|z_1|=|z_2|$, consider
\[
 p|_{[0,+\infty)}:[0,+\infty)\longrightarrow\mathbb C/{\sim}.
\]
\noindent\textbf{(1)} For $x,y\in[0,+\infty)$,
$x\ne y\Longrightarrow x\not\sim y$.

\noindent\textbf{(2)} For every $z\in\mathbb C$ there is
$t\in[0,+\infty)$ with $z\sim t$: take $t=|z|$.
\[
\begin{tikzcd}[column sep=large,row sep=small]
 \mathbb C\arrow[r,"p"]\arrow[d,"{f(z)=|z|}"']
    & \mathbb C/{\sim} \\
 {[0,+\infty)}\arrow[ur,dashed,"{p|_{[0,+\infty)}}"'] &
\end{tikzcd}
\]
If $|z_1|=|z_2|$, then $f(z_1)=f(z_2)$.

\medskip
\noindent\textbf{(iv) Complex numbers with equal argument.}
Put
\[
 S'=\{z\in\mathbb C\mid |z|=1\},\qquad
 z_1\sim z_2\iff\arg(z_1)=\arg(z_2)
 \quad(z_1,z_2\in\mathbb C^*).
\]
\noindent\textbf{(1)} If $z_1,z_2\in S'$ and $z_1\ne z_2$,
then $z_1\not\sim z_2$ (their arguments differ).

\noindent\textbf{(2)} For every $z\in\mathbb C^*$ there exists
$z_1\in S'$ with $z\sim z_1$. Indeed, write
\[
 z=r(\cos t+i\sin t),\qquad z_1=\cos t+i\sin t.
\]
Thus the notes identify $\mathbb C^*/{\sim}$ with $S'$ via
representatives on the unit circle:
\[
\begin{tikzcd}[column sep=large,row sep=small]
 \mathbb C^*\arrow[r,"p"]\arrow[d,"{f(z)=z/|z|}"']
    & \mathbb C^*/{\sim} \\
 S'\arrow[ur,dashed,"{p|_{S'}}"'] &
\end{tikzcd}
\]

\endgroup
; no preamble or forced page break.
% EDITORIAL NOTES:
% - This page continues examples (i), (iii), and (iv) of the preceding page.
% - In the opening line, square brackets and braces denote the integer and
%   fractional parts; the compact handwritten expression is interpreted as
%   [x]+{x}-[y]-{y} being an integer when x~y.
% - The proof explicitly shows reflexivity and symmetry of congruence mod Z;
%   transitivity is not written out on this manuscript page.
% - A small marginal reminder repeats (iv), x~y iff x-y in Q, and R/~ ~ R.
%   It is not reproduced as another numbered example.
% - The diagrams are redrawn for legibility; p restricted to a chosen set
%   of representatives is the indicated bijection.
% - Equality of complex arguments depends on a consistent choice of argument
%   (or equality modulo 2pi); the notes leave this convention implicit.
% - In the last diagram the handwritten notation for the map to S' is faint;
%   the unit-normalisation map z -> z/|z| is consistent with the proof.

\subsection*{Seminar 3 - Page 6}
\begingroup
\setlength{\abovedisplayskip}{5pt}
\setlength{\belowdisplayskip}{5pt}
\setlength{\abovedisplayshortskip}{3pt}
\setlength{\belowdisplayshortskip}{3pt}

\noindent\textbf{Proof (continuation of example (i)).}
Write $x=\lfloor x\rfloor+\{x\}$ and $y=\lfloor y\rfloor+\{y\}$.
If $x-y\in\mathbb Z$, then
\[
 \lfloor x\rfloor+\{x\}-\lfloor y\rfloor-\{y\}\in\mathbb Z,
 \qquad \{x\}-\{y\}\in\mathbb Z.
\]
For the relation $x\sim y\iff x-y\in\mathbb Z$:
\[
 \begin{aligned}
 \text{(R)}\quad &x\sim x &&(x-x=0\in\mathbb Z),\\
 \text{(S)}\quad &x\sim y\Longrightarrow y\sim x
     &&(x-y\in\mathbb Z\Longrightarrow y-x\in\mathbb Z).
 \end{aligned}
\]

\noindent\underline{Variant 1.} The interval $[0,1)$ is a system of
representatives for $\sim$ if and only if
\[
 p|_{[0,1)}:[0,1)\longrightarrow\mathbb R/{\sim}
 \quad\text{is bijective}.
\]
\noindent\textbf{Step 1.} If $x,y\in[0,1)$ and $x\ne y$, then
$x\not\sim y$ (evident).

\noindent\textbf{Step 2.} For every $x\in\mathbb R$ there is
$t\in[0,1)$ with $x-t\in\mathbb Z$: take $t=\{x\}$.

\medskip
\noindent\textbf{(iii) Complex numbers with equal modulus.}
For the quotient $\mathbb C/{\sim}$, where
$z_1\sim z_2\iff|z_1|=|z_2|$, consider
\[
 p|_{[0,+\infty)}:[0,+\infty)\longrightarrow\mathbb C/{\sim}.
\]
\noindent\textbf{(1)} For $x,y\in[0,+\infty)$,
$x\ne y\Longrightarrow x\not\sim y$.

\noindent\textbf{(2)} For every $z\in\mathbb C$ there is
$t\in[0,+\infty)$ with $z\sim t$: take $t=|z|$.
\[
\begin{tikzcd}[column sep=large,row sep=small]
 \mathbb C\arrow[r,"p"]\arrow[d,"{f(z)=|z|}"']
    & \mathbb C/{\sim} \\
 {[0,+\infty)}\arrow[ur,dashed,"{p|_{[0,+\infty)}}"'] &
\end{tikzcd}
\]
If $|z_1|=|z_2|$, then $f(z_1)=f(z_2)$.

\medskip
\noindent\textbf{(iv) Complex numbers with equal argument.}
Put
\[
 S'=\{z\in\mathbb C\mid |z|=1\},\qquad
 z_1\sim z_2\iff\arg(z_1)=\arg(z_2)
 \quad(z_1,z_2\in\mathbb C^*).
\]
\noindent\textbf{(1)} If $z_1,z_2\in S'$ and $z_1\ne z_2$,
then $z_1\not\sim z_2$ (their arguments differ).

\noindent\textbf{(2)} For every $z\in\mathbb C^*$ there exists
$z_1\in S'$ with $z\sim z_1$. Indeed, write
\[
 z=r(\cos t+i\sin t),\qquad z_1=\cos t+i\sin t.
\]
Thus the notes identify $\mathbb C^*/{\sim}$ with $S'$ via
representatives on the unit circle:
\[
\begin{tikzcd}[column sep=large,row sep=small]
 \mathbb C^*\arrow[r,"p"]\arrow[d,"{f(z)=z/|z|}"']
    & \mathbb C^*/{\sim} \\
 S'\arrow[ur,dashed,"{p|_{S'}}"'] &
\end{tikzcd}
\]

\endgroup
; no preamble or forced page break.
% EDITORIAL NOTES:
% - This page continues examples (i), (iii), and (iv) of the preceding page.
% - In the opening line, square brackets and braces denote the integer and
%   fractional parts; the compact handwritten expression is interpreted as
%   [x]+{x}-[y]-{y} being an integer when x~y.
% - The proof explicitly shows reflexivity and symmetry of congruence mod Z;
%   transitivity is not written out on this manuscript page.
% - A small marginal reminder repeats (iv), x~y iff x-y in Q, and R/~ ~ R.
%   It is not reproduced as another numbered example.
% - The diagrams are redrawn for legibility; p restricted to a chosen set
%   of representatives is the indicated bijection.
% - Equality of complex arguments depends on a consistent choice of argument
%   (or equality modulo 2pi); the notes leave this convention implicit.
% - In the last diagram the handwritten notation for the map to S' is faint;
%   the unit-normalisation map z -> z/|z| is consistent with the proof.

\subsection*{Seminar 3 - Page 6}
\begingroup
\setlength{\abovedisplayskip}{5pt}
\setlength{\belowdisplayskip}{5pt}
\setlength{\abovedisplayshortskip}{3pt}
\setlength{\belowdisplayshortskip}{3pt}

\noindent\textbf{Proof (continuation of example (i)).}
Write $x=\lfloor x\rfloor+\{x\}$ and $y=\lfloor y\rfloor+\{y\}$.
If $x-y\in\mathbb Z$, then
\[
 \lfloor x\rfloor+\{x\}-\lfloor y\rfloor-\{y\}\in\mathbb Z,
 \qquad \{x\}-\{y\}\in\mathbb Z.
\]
For the relation $x\sim y\iff x-y\in\mathbb Z$:
\[
 \begin{aligned}
 \text{(R)}\quad &x\sim x &&(x-x=0\in\mathbb Z),\\
 \text{(S)}\quad &x\sim y\Longrightarrow y\sim x
     &&(x-y\in\mathbb Z\Longrightarrow y-x\in\mathbb Z).
 \end{aligned}
\]

\noindent\underline{Variant 1.} The interval $[0,1)$ is a system of
representatives for $\sim$ if and only if
\[
 p|_{[0,1)}:[0,1)\longrightarrow\mathbb R/{\sim}
 \quad\text{is bijective}.
\]
\noindent\textbf{Step 1.} If $x,y\in[0,1)$ and $x\ne y$, then
$x\not\sim y$ (evident).

\noindent\textbf{Step 2.} For every $x\in\mathbb R$ there is
$t\in[0,1)$ with $x-t\in\mathbb Z$: take $t=\{x\}$.

\medskip
\noindent\textbf{(iii) Complex numbers with equal modulus.}
For the quotient $\mathbb C/{\sim}$, where
$z_1\sim z_2\iff|z_1|=|z_2|$, consider
\[
 p|_{[0,+\infty)}:[0,+\infty)\longrightarrow\mathbb C/{\sim}.
\]
\noindent\textbf{(1)} For $x,y\in[0,+\infty)$,
$x\ne y\Longrightarrow x\not\sim y$.

\noindent\textbf{(2)} For every $z\in\mathbb C$ there is
$t\in[0,+\infty)$ with $z\sim t$: take $t=|z|$.
\[
\begin{tikzcd}[column sep=large,row sep=small]
 \mathbb C\arrow[r,"p"]\arrow[d,"{f(z)=|z|}"']
    & \mathbb C/{\sim} \\
 {[0,+\infty)}\arrow[ur,dashed,"{p|_{[0,+\infty)}}"'] &
\end{tikzcd}
\]
If $|z_1|=|z_2|$, then $f(z_1)=f(z_2)$.

\medskip
\noindent\textbf{(iv) Complex numbers with equal argument.}
Put
\[
 S'=\{z\in\mathbb C\mid |z|=1\},\qquad
 z_1\sim z_2\iff\arg(z_1)=\arg(z_2)
 \quad(z_1,z_2\in\mathbb C^*).
\]
\noindent\textbf{(1)} If $z_1,z_2\in S'$ and $z_1\ne z_2$,
then $z_1\not\sim z_2$ (their arguments differ).

\noindent\textbf{(2)} For every $z\in\mathbb C^*$ there exists
$z_1\in S'$ with $z\sim z_1$. Indeed, write
\[
 z=r(\cos t+i\sin t),\qquad z_1=\cos t+i\sin t.
\]
Thus the notes identify $\mathbb C^*/{\sim}$ with $S'$ via
representatives on the unit circle:
\[
\begin{tikzcd}[column sep=large,row sep=small]
 \mathbb C^*\arrow[r,"p"]\arrow[d,"{f(z)=z/|z|}"']
    & \mathbb C^*/{\sim} \\
 S'\arrow[ur,dashed,"{p|_{S'}}"'] &
\end{tikzcd}
\]

\endgroup
; no preamble or forced page break.
% EDITORIAL NOTES:
% - This page continues examples (i), (iii), and (iv) of the preceding page.
% - In the opening line, square brackets and braces denote the integer and
%   fractional parts; the compact handwritten expression is interpreted as
%   [x]+{x}-[y]-{y} being an integer when x~y.
% - The proof explicitly shows reflexivity and symmetry of congruence mod Z;
%   transitivity is not written out on this manuscript page.
% - A small marginal reminder repeats (iv), x~y iff x-y in Q, and R/~ ~ R.
%   It is not reproduced as another numbered example.
% - The diagrams are redrawn for legibility; p restricted to a chosen set
%   of representatives is the indicated bijection.
% - Equality of complex arguments depends on a consistent choice of argument
%   (or equality modulo 2pi); the notes leave this convention implicit.
% - In the last diagram the handwritten notation for the map to S' is faint;
%   the unit-normalisation map z -> z/|z| is consistent with the proof.

\subsection*{Seminar 3 - Page 6}
\begingroup
\setlength{\abovedisplayskip}{5pt}
\setlength{\belowdisplayskip}{5pt}
\setlength{\abovedisplayshortskip}{3pt}
\setlength{\belowdisplayshortskip}{3pt}

\noindent\textbf{Proof (continuation of example (i)).}
Write $x=\lfloor x\rfloor+\{x\}$ and $y=\lfloor y\rfloor+\{y\}$.
If $x-y\in\mathbb Z$, then
\[
 \lfloor x\rfloor+\{x\}-\lfloor y\rfloor-\{y\}\in\mathbb Z,
 \qquad \{x\}-\{y\}\in\mathbb Z.
\]
For the relation $x\sim y\iff x-y\in\mathbb Z$:
\[
 \begin{aligned}
 \text{(R)}\quad &x\sim x &&(x-x=0\in\mathbb Z),\\
 \text{(S)}\quad &x\sim y\Longrightarrow y\sim x
     &&(x-y\in\mathbb Z\Longrightarrow y-x\in\mathbb Z).
 \end{aligned}
\]

\noindent\underline{Variant 1.} The interval $[0,1)$ is a system of
representatives for $\sim$ if and only if
\[
 p|_{[0,1)}:[0,1)\longrightarrow\mathbb R/{\sim}
 \quad\text{is bijective}.
\]
\noindent\textbf{Step 1.} If $x,y\in[0,1)$ and $x\ne y$, then
$x\not\sim y$ (evident).

\noindent\textbf{Step 2.} For every $x\in\mathbb R$ there is
$t\in[0,1)$ with $x-t\in\mathbb Z$: take $t=\{x\}$.

\medskip
\noindent\textbf{(iii) Complex numbers with equal modulus.}
For the quotient $\mathbb C/{\sim}$, where
$z_1\sim z_2\iff|z_1|=|z_2|$, consider
\[
 p|_{[0,+\infty)}:[0,+\infty)\longrightarrow\mathbb C/{\sim}.
\]
\noindent\textbf{(1)} For $x,y\in[0,+\infty)$,
$x\ne y\Longrightarrow x\not\sim y$.

\noindent\textbf{(2)} For every $z\in\mathbb C$ there is
$t\in[0,+\infty)$ with $z\sim t$: take $t=|z|$.
\[
\begin{tikzcd}[column sep=large,row sep=small]
 \mathbb C\arrow[r,"p"]\arrow[d,"{f(z)=|z|}"']
    & \mathbb C/{\sim} \\
 {[0,+\infty)}\arrow[ur,dashed,"{p|_{[0,+\infty)}}"'] &
\end{tikzcd}
\]
If $|z_1|=|z_2|$, then $f(z_1)=f(z_2)$.

\medskip
\noindent\textbf{(iv) Complex numbers with equal argument.}
Put
\[
 S'=\{z\in\mathbb C\mid |z|=1\},\qquad
 z_1\sim z_2\iff\arg(z_1)=\arg(z_2)
 \quad(z_1,z_2\in\mathbb C^*).
\]
\noindent\textbf{(1)} If $z_1,z_2\in S'$ and $z_1\ne z_2$,
then $z_1\not\sim z_2$ (their arguments differ).

\noindent\textbf{(2)} For every $z\in\mathbb C^*$ there exists
$z_1\in S'$ with $z\sim z_1$. Indeed, write
\[
 z=r(\cos t+i\sin t),\qquad z_1=\cos t+i\sin t.
\]
Thus the notes identify $\mathbb C^*/{\sim}$ with $S'$ via
representatives on the unit circle:
\[
\begin{tikzcd}[column sep=large,row sep=small]
 \mathbb C^*\arrow[r,"p"]\arrow[d,"{f(z)=z/|z|}"']
    & \mathbb C^*/{\sim} \\
 S'\arrow[ur,dashed,"{p|_{S'}}"'] &
\end{tikzcd}
\]

\endgroup
